\section{A product launch changes the history that must be explained}
\label{sec:apple-competition-transition}

Keeping a supply route open does not ensure that it supplies the product
customers now want. A new offering can replace demand for an earlier one while
the company still holds its components, finished goods or purchase commitments.
Apple's product-transition risk explicitly says that new products, services
and technologies may replace existing offerings and produce lower revenues
and lower profit margins \citep[Risk Factors, printed p.~7]{apple2024}.
A model that equates a successful release with an improvement in financial
performance would remove a possibility the source asks it to preserve.

The preceding competitive-market discussion explains why a transition is not
merely a production date. Apple describes price competition, short product
life cycles, changing standards, rapid adoption of technology and price-sensitive
customers. It says that developing almost the whole product solution requires
substantial research and development, without assurance of the expected returns
or successful marketing. Its reported intellectual-property portfolio also
does not assure protection in every jurisdiction. The discussion of competitors
imitating products or infringing rights remains an issuer assertion; it does
not decide a particular infringement dispute. Apple's description of minority
share in three global device markets and of competitors able to sell at little
or no profit supplies competitive context, rather than a numerical market
share or a forecast of Apple's margins. Services face a further set of
alternatives, including free content. The company's expectation of more
intense competition is a dated expectation, not an observed future result.

The transition passage identifies several separate conditions for a new
offering: development must succeed on time, customers must accept it, production
must reach suitable quantities and costs, and application software or other
third-party support must be available. New technology can create ramp-up or
quality problems. Purchase commitments and inventory must also be managed in
line with anticipated demand. Note 12 makes the supply condition concrete:
new products often use custom components available from only one source;
capacity can be constrained while yields mature or manufacturing expands,
and a supplier may instead concentrate on common components
\citep[Note 12, printed p.~46]{apple2024}. A forecast of customer acceptance
cannot establish that those components will be available.

Suppose, as an illustration, that customers switch from an earlier product to
its replacement. If existing components can be used in the replacement and
the supplier can expand output, the transition presents one set of feasible
production choices. If the components are incompatible or the new component
has constrained capacity, the same assumed customer response presents another.
Neither history is identified by aggregate revenue or inventory. The comparison
requires the old and new product definitions, component compatibility, inventory
cohorts, committed purchases, cancellation terms, lead times and available
capacity. Prices, quantities and costs would then be separate scenario inputs
or model consequences. The risk language gives no substitution fraction,
write-down amount or probability with which to combine the histories.

Even a delivered replacement has more than one financial date. Apple's revenue
policy generally connects product revenue to transfer of control, while
allocated revenue for bundled services is deferred over the estimated service
period; collection follows its own terms
\citep[Note 2, printed p.~35]{apple2024}. Development, announcement, production,
shipment, service provision and cash receipt therefore cannot share a single
``launch'' timestamp. The reported policy must be applied to the relevant
contract and performance obligations. Any inventory disposition, commitment
change or development-cost treatment needs its own facts and applicable
accounting method. Preserving these distinctions lets the model ask whether
an assumed transition is feasible and how its financial consequences would
arise, while leaving actual product histories unresolved.
